\chapter{Making Decisions}\label{ch:conditions}

\begin{goals}
\begin{itemize}
  \item Running instructions only when a condition holds, with \code{if}
  \item Choosing between two paths with \code{else}, and among several with
        \code{else condition}
  \item Nested conditions and combined conditions
  \item Choosing a value with the \code{? :} operator
  \item Choosing among fixed cases with \code{match}
\end{itemize}
\end{goals}

\section{A program that decides}

Every program we have written so far had only one path: it started at the
first line and went line by line to the end. Real programs, however, have to
make decisions. ``If the grade is ten or more, print pass.'' ``If the
password is right, log in.'' ``If the balance is not enough, give a
warning.'' These ``ifs'' are exactly the word Salam uses.

\section{if}

The simplest form of decision is this: if a condition holds, run one or more
instructions; otherwise do nothing:

\program{The simplest if}
\begin{salamcode}
func main:
    grade := 14
    if grade >= 10:
        println "Pass"
    end
    println "End of program"
end
\end{salamcode}

\begin{salamout}
Pass
End of program
\end{salamout}

The structure of \code{if} is the same pattern you saw in \code{func main}:
a line that ends with a colon, some indented instructions, and \code{end}.
What comes between \code{if} and the colon is the \textbf{condition}: a
boolean expression that is either true or false. If it is true, the
instructions inside the block run; if it is false, Salam skips the whole
block and carries on from the line after \code{end}.

Change the grade in the program above to \code{8} and run it again. This
time only ``End of program'' is printed, because the condition did not hold
and the instruction \code{println "Pass"} never ran at all.

\begin{note}
The condition of an \code{if} can be anything that gives a boolean value: a
comparison such as \code{grade >= 10}, a boolean variable such as
\code{has ticket}, or a combination of these with \code{and} and
\code{or}. A number on its own is not a condition; writing \code{if grade:}
is an error.
\end{note}

\section{if and else}

Many decisions have two paths: either this, or that. The word \kwd{else}
marks the second path:

\program{if and else}
\begin{salamcode}
func main:
    grade := 8
    if grade >= 10:
        println "Pass"
    else:
        println "Fail"
    end
end
\end{salamcode}

\begin{salamout}
Fail
\end{salamout}

Notice that \code{else} has a colon of its own but no separate \code{end};
a single \code{end} at the bottom closes the whole ``if, else'' structure.
Of the two blocks, exactly one always runs: never both, never neither.

\program{Even or odd}
\begin{salamcode}
func main:
    number := 27
    if (number % 2) == 0:
        println number, "is even"
    else:
        println number, "is odd"
    end
end
\end{salamcode}

\begin{salamout}
27 is odd
\end{salamout}

An even number is one whose remainder on division by 2 is zero. This is one
of the most frequently used conditions in programming, and you will see it
many times.

\section{More than two paths: else with a condition}

Sometimes there are more than two paths. Say we want to sort a grade into
``excellent'', ``good'', ``acceptable'' and ``fail''. A new condition can be
written after \code{else}:

\program{A chain of conditions}
\begin{salamcode}
func main:
    grade := 17
    if grade >= 18:
        println "Excellent"
    else grade >= 15:
        println "Good"
    else grade >= 10:
        println "Acceptable"
    else:
        println "Fail"
    end
end
\end{salamcode}

\begin{salamout}
Good
\end{salamout}

Salam tests the conditions from top to bottom and settles for the
\emph{first} one that is true; the rest are not tested at all. A grade of 17
is less than 18, so the first condition is false; but it is more than 15, so
``Good'' is printed and Salam passes over the remaining conditions. The
final \code{else}, which has no condition, runs when none of the conditions
holds. Writing it is optional.

\begin{warn}
The order of the conditions in a chain matters. Had we started the program
above with the condition \code{grade >= 10}, a grade of 17 would also have
been ``Acceptable'' and the later conditions would never have been reached.
Write the conditions from the most specific to the most general.
\end{warn}

\section{if on one line}

When the block of an \code{if} holds only one short instruction, the whole
thing can be written on one line. The \code{end} is still required:

\program{A one-line if}
\begin{salamcode}
func main:
    temp := 35
    if temp > 30: println "It is hot" end
    if temp < 10: println "It is cold" end
    if (temp >= 10) and (temp <= 30): println "It is mild" end
end
\end{salamcode}

\begin{salamout}
It is hot
\end{salamout}

This form suits short, simple conditions. As soon as the block has more than
one instruction or the condition grows long, go back to the multi-line
form.

\section{A condition inside a condition}

Any instruction can be written inside the block of an \code{if}, including
another \code{if}. This is called \textbf{nesting}:

\program{A nested condition}
\begin{salamcode}
func main:
    age := 20
    has ticket := true
    if age >= 18:
        if has ticket:
            println "Welcome"
        else:
            println "Please buy a ticket"
        end
    else:
        println "No entry under 18"
    end
end
\end{salamcode}

\begin{salamout}
Welcome
\end{salamout}

Look at the indentation: the inner \code{if} is one step further in, and
the \code{end} of each block sits exactly below its own \code{if}. Without
proper indentation, working out which \code{end} closes which block becomes
almost impossible.

Many nested conditions can be flattened with \code{and} and \code{or}. The
same program, without nesting:

\begin{salamcode}
func main:
    age := 20
    has ticket := true
    if (age >= 18) and has ticket:
        println "Welcome"
    else age < 18:
        println "No entry under 18"
    else:
        println "Please buy a ticket"
    end
end
\end{salamcode}

\begin{salamout}
Welcome
\end{salamout}

Neither form is ``more correct''. A rule of thumb: write whichever one you
understand better on a single reading.

\section{The largest of several numbers}

A common pattern built from \code{if} and a changeable variable: finding the
largest value. First we assume the first one is the largest, then we weigh
the others against it:

\program{The largest of three numbers}
\begin{salamcode}
func main:
    a := 12
    b := 25
    c := 19
    mut largest := a
    if b > largest: largest = b end
    if c > largest: largest = c end
    println "Largest:", largest
end
\end{salamcode}

\begin{salamout}
Largest: 25
\end{salamout}

Try it with any values for \code{a}, \code{b} and \code{c}. This pattern
will serve us again in Chapters~\ref{ch:loops} and~\ref{ch:arrays}, when
instead of three numbers we have hundreds.

\section{Choosing a value: the \code{?} operator}

Sometimes we do not want an instruction to run; we want to choose one of two
\emph{values} according to a condition. In the last chapter you saw the
operator \code{condition ? a : b}, which does exactly this:

\program{The choice operator}
\begin{salamcode}
func main:
    age := 15
    ticket price := age < 12 ? 5 : 12
    println "Ticket price:", ticket price
    println age >= 18 ? "adult" : "teenager"
end
\end{salamcode}

\begin{salamout}
Ticket price: 12
teenager
\end{salamout}

Read the line \code{ticket price := age < 12 ? 5 : 12} like this: ``the
ticket price is: if the age is under 12, then 5, otherwise 12.'' The same
could have been done with an \code{if}, an \code{else} and a changeable
variable, but this form is shorter.

\begin{tip}
Use the \code{? :} operator only for simple choices. As soon as two or three
of them are nested inside one another, go back to \code{if}; readability
matters more than brevity.
\end{tip}

\section{match: choosing among fixed cases}

When one value is compared against several \emph{fixed} cases (the number of
a day of the week, the number of a month, the code of a menu option), a
chain of \code{if} and \code{else} grows long and monotonous. For this
Salam has the word \kwd{match}:

\program{The name of a day of the week}
\begin{salamcode}
func main:
    day := 5
    match day:
        1: println "Monday" end
        2: println "Tuesday" end
        3: println "Wednesday" end
        4: println "Thursday" end
        5: println "Friday" end
        6: println "Saturday" end
        7: println "Sunday" end
        else: println "Invalid day number" end
    end
end
\end{salamcode}

\begin{salamout}
Friday
\end{salamout}

\code{match} compares the value of \code{day} with each of the cases and
runs the block of the case that is equal. Each case is a small block that
starts with a colon and is closed with \code{end}; it can be written on one
line or on several. \code{else} is for when no case is equal.

\code{match} can also be used where a \emph{value} is expected, instead of
running an instruction. The cases are written exactly the same way; the value
of a case is the expression its block gives back. Several values can share one
case, separated by commas:

\program{The season of each month}
\begin{salamcode}
func main:
    month := 4
    season := match month:
        3, 4, 5: "spring" end
        6, 7, 8: "summer" end
        9, 10, 11: "autumn" end
        12, 1, 2: "winter" end
        else: "invalid" end
    end
    println "Month", month, "is in", season
end
\end{salamcode}

\begin{salamout}
Month 4 is in spring
\end{salamout}

This form is very convenient for ``conversion tables'' (number to name, code
to message). Written with \code{if}, such a table would need twelve
conditions and twelve \code{println} instructions.

\section{A more complete program}

The body mass index is found by dividing the weight (in kilograms) by the
square of the height (in metres). This program works out the index and
classifies it:

\program{Body mass index}
\begin{salamcode}
func main:
    weight := 70.0
    height := 1.75
    index := weight / (height * height)
    println "Body mass index:", index

    if index < 18.5:
        println "Category: underweight"
    else index < 25.0:
        println "Category: normal"
    else index < 30.0:
        println "Category: overweight"
    else:
        println "Category: obese"
    end
end
\end{salamcode}

\begin{salamout}
Body mass index: 22.857142857142858
Category: normal
\end{salamout}

Notice that in the chain of conditions there was no need to write
\begin{center}\code{index >= 18.5 and index < 25.0}\end{center}
because if we have reached the second condition, the first one was false,
and the index is certainly not below 18.5. The chain guarantees this by
itself.

\begin{summary}
\begin{itemize}
  \item \code{if condition:} runs a block only when the condition is true.
        \code{else:} gives a second path and \code{else condition:} gives
        further paths. One \code{end} closes the whole structure.
  \item In a chain, the first true condition wins; the order matters.
  \item Short blocks can be written on one line, with \code{end} at the
        end.
  \item An \code{if} can sit inside another \code{if}; keep the
        indentation right. Many nestings can be simplified with
        \code{and} and \code{or}.
  \item \code{condition ? a : b} chooses one of two values.
  \item \code{match} compares one value against several fixed cases; each
        case is a block that opens with a colon and closes with \code{end},
        and it serves both as an instruction and as a value.
\end{itemize}
\end{summary}

\section*{Exercises}
\markright{Exercises}

\exercise Put a number into a variable and print whether it is positive,
negative or zero.

\exercise Put the three sides of a triangle into three variables and print
whether the triangle is equilateral (all three sides equal), isosceles (two
sides equal) or scalene.

\exercise Put the hour of the day (0 to 23) into a variable and print the
right greeting: before 12 ``Good morning'', from 12 until before 17 ``Good
afternoon'', from 17 until before 20 ``Good evening'', and from 20 onwards
``Good night''.

\exercise With \code{match}, convert the number of a month into its number
of days: 30 days for April, June, September and November, 28 for February,
and 31 for the rest.

\exercise Put the weight of a parcel into a variable and work out the
postage: up to 1 kilogram costs 3, up to 5 kilograms costs 5, and above that
it costs 5 plus 2 for every kilogram over 5. If the parcel is going to
another city (a boolean variable), add 20 percent to the cost.

\exercise What does the program below print? Guess first, then run it, and
if your guess was wrong, find out why:
\begin{salamcode}
func main:
    number := 15
    if number > 10:
        println "greater than ten"
    else number > 5:
        println "greater than five"
    end
    if (number % 5) == 0: println "a multiple of five" end
    if (number % 3) == 0: println "a multiple of three" end
end
\end{salamcode}
